% !TEX TS-program = XeLaTeX
% use the following command:
% all document files must be coded in UTF-8
\documentclass[spanish]{textolivre}
% build HTML with: make4ht -e build.lua -c textolivre.cfg -x -u article "fn-in,svg,pic-align"

\journalname{Texto Livre}
\thevolume{19}
%\thenumber{1} % old template
\theyear{2026}
\receiveddate{\DTMdisplaydate{2025}{9}{16}{-1}} % YYYY MM DD
\accepteddate{\DTMdisplaydate{2026}{1}{6}{-1}}
\publisheddate{\DTMdisplaydate{2026}{4}{25}{-1}}
\corrauthor{Andrés Seguel-Arriagada}
\articledoi{10.1590/1983-3652.2026.61792}
%\articleid{NNNN} % if the article ID is not the last 5 numbers of its DOI, provide it using \articleid{} commmand 
% list of available sesscions in the journal: articles, dossier, reports, essays, reviews, interviews, editorial
\articlesessionname{essays}
\runningauthor{Seguel-Arriagada} 
%\editorname{Leonardo Araújo} % old template
\sectioneditorname{Hugo Heredia Ponce~\orcid{0000-0003-3657-1369}}
\layouteditorname{Daniervelin Pereira~\orcid{0000-0003-1861-3609}}

\title{La Inteligencia Artificial Generativa en la Sociedad 5.0: hacia una humanidad ampliada para un profesional sensible}
\othertitle{A Inteligência Artificial Generativa na Sociedade 5.0: rumo a uma humanidade ampliada para um profissional sensível}
\othertitle{Generative Artificial Intelligence in Society 5.0: toward an augmented humanity for a sensitive professional}
% if there is a third language title, add here:
%\othertitle{Artikelvorlage zur Einreichung beim Texto Livre Journal}

\author[1]{Andrés Seguel-Arriagada~\orcid{0000-0003-2549-5890 }\thanks{Email: \href{mailto:andres.seguel@uvm.cl}{andres.seguel@uvm.cl}}}
\affil[1]{Universidad Viña del Mar, Viña del Mar, Valparaíso, Chile.}

\addbibresource{article.bib}
% use biber instead of bibtex
% $ biber article

% used to create dummy text for the template file
\definecolor{dark-gray}{gray}{0.35} % color used to display dummy texts
\usepackage{lipsum}
\SetLipsumParListSurrounders{\colorlet{oldcolor}{.}\color{dark-gray}}{\color{oldcolor}}

% used here only to provide the XeLaTeX and BibTeX logos
\usepackage{hologo}

% if you use multirows in a table, include the multirow package
\usepackage{multirow}

% provides sidewaysfigure environment
\usepackage{rotating}

% CUSTOM EPIGRAPH - BEGIN 
%%% https://tex.stackexchange.com/questions/193178/specific-epigraph-style
\usepackage{epigraph}
\renewcommand\textflush{flushright}
\makeatletter
\newlength\epitextskip
\pretocmd{\@epitext}{\em}{}{}
\apptocmd{\@epitext}{\em}{}{}
\patchcmd{\epigraph}{\@epitext{#1}\\}{\@epitext{#1}\\[\epitextskip]}{}{}
\makeatother
\setlength\epigraphrule{0pt}
\setlength\epitextskip{0.5ex}
\setlength\epigraphwidth{.7\textwidth}
% CUSTOM EPIGRAPH - END

% to use IPA symbols in unicode add
%\usepackage{fontspec}
%\newfontfamily\ipafont{CMU Serif}
%\newcommand{\ipa}[1]{{\ipafont #1}}
% and in the text you may use the \ipa{...} command passing the symbols in unicode

% LANGUAGE - BEGIN
% ARABIC
% for languages that use special fonts, you must provide the typeface that will be used
% \setotherlanguage{arabic}
% \newfontfamily\arabicfont[Script=Arabic]{Amiri}
% \newfontfamily\arabicfontsf[Script=Arabic]{Amiri}
% \newfontfamily\arabicfonttt[Script=Arabic]{Amiri}
%
% in the article, to add arabic text use: \textlang{arabic}{ ... }
%
% RUSSIAN
% for russian text we also need to define fonts with support for Cyrillic script
% \usepackage{fontspec}
% \setotherlanguage{russian}
% \newfontfamily\cyrillicfont{Times New Roman}
% \newfontfamily\cyrillicfontsf{Times New Roman}[Script=Cyrillic]
% \newfontfamily\cyrillicfonttt{Times New Roman}[Script=Cyrillic]
%
% in the text use \begin{russian} ... \end{russian}
% LANGUAGE - END

% EMOJIS - BEGIN
% to use emoticons in your manuscript
% https://stackoverflow.com/questions/190145/how-to-insert-emoticons-in-latex/57076064
% using font Symbola, which has full support
% the font may be downloaded at:
% https://dn-works.com/ufas/
% add to preamble:
% \newfontfamily\Symbola{Symbola}
% in the text use:
% {\Symbola }
% EMOJIS - END

% LABEL REFERENCE TO DESCRIPTIVE LIST - BEGIN
% reference itens in a descriptive list using their labels instead of numbers
% insert the code below in the preambule:
%\makeatletter
%\let\orgdescriptionlabel\descriptionlabel
%\renewcommand*{\descriptionlabel}[1]{%
%  \let\orglabel\label
%  \let\label\@gobble
%  \phantomsection
%  \edef\@currentlabel{#1\unskip}%
%  \let\label\orglabel
%  \orgdescriptionlabel{#1}%
%}
%\makeatother
%
% in your document, use as illustraded here:
%\begin{description}
%  \item[first\label{itm1}] this is only an example;
%  % ...  add more items
%\end{description}
% LABEL REFERENCE TO DESCRIPTIVE LIST - END


% add line numbers for submission
%\usepackage{lineno}
%\linenumbers

\begin{document}
\maketitle

\begin{polyabstract}
\begin{abstract}
En el marco de la Sociedad 5.0, caracterizada por la integración de tecnologías disruptivas orientadas al bienestar humano, este ensayo propone una reflexión crítica y sensible sobre el lugar de la subjetividad en la era de la Inteligencia Artificial Generativa (IAG). Lejos de limitarse a lo técnico, se argumenta que la irrupción de la IAG interpela profundamente nuestras nociones de sujeto, alteridad y encuentro. Al partir del enfoque filosófico y ético, el texto desarrolla la idea de humanidad ampliada como una expansión de la sensibilidad, la responsabilidad y la co-existencia con inteligencias no-humanas. Desde los aportes de Merleau-Ponty, Levinas y Maturana, se propone una pedagogía de la sensibilidad como vía para formar profesionales sensibles, capaces de sostener vínculos éticos y significativos en entornos híbridos. Frente al riesgo de una tecnocracia deshumanizante, se plantea que la verdadera revolución pendiente no es la digital, sino la del corazón: una apertura al otro – humano o artificial – como ocasión de subjetivación, cuidado y creación compartida. El ensayo concluye que la IAG, bien habitada, puede ser un espejo y un puente hacia formas inéditas de humanidad.

\keywords{Humanidad \sep IAG \sep Pedagogía \sep Subjetividad \sep Sociedad}
\end{abstract}

\begin{portuguese}
\begin{abstract}
No âmbito da Sociedade 5.0, caracterizada pela integração de tecnologias disruptivas voltadas para o bem-estar humano, este ensaio propõe uma reflexão crítica e sensível sobre o lugar da subjetividade na era da Inteligência Artificial Generativa (IAG). Longe de se limitar aos aspectos técnicos, argumenta que a emergência da IAG desafia profundamente nossas noções de sujeito, alteridade e encontro. De uma perspectiva filosófica e ética, o texto desenvolve a ideia de humanidade expandida como uma expansão da sensibilidade, da responsabilidade e da coexistência com inteligências não humanas. Inspirando-se nas contribuições de Merleau-Ponty, Levinas e Maturana, propõe uma pedagogia da sensibilidade como meio de formar profissionais sensíveis, capazes de sustentar relações éticas e significativas em ambientes híbridos. Diante do risco de uma tecnocracia desumanizadora, postula que a verdadeira revolução por vir não é digital, mas sim uma revolução do coração: uma abertura ao outro — humano ou artificial — como oportunidade para subjetivação, cuidado e criação compartilhada. O ensaio conclui que a IAG, quando bem habitada, pode ser um espelho e uma ponte para formas de humanidade sem precedentes.

\keywords{Humanidade \sep IAG \sep Pedagogia \sep Subjetividade \sep Sociedade}
\end{abstract}
\end{portuguese}

\begin{english}
\begin{abstract}
Within the framework of Society 5.0, characterized by the integration of disruptive technologies oriented towards human well-being, this essay proposes a critical and sensitive reflection on the place of subjectivity in the era of Generative Artificial Intelligence. Far from being limited to the technical, it argues that the irruption of GGI deeply challenges our notions of subject, otherness and encounter. Starting from a philosophical and ethical approach, the text develops the idea of expanded humanity as an expansion of sensitivity, responsibility and co-existence with non-human intelligences. From the contributions of Merleau-Ponty, Levinas and Maturana, a pedagogy of sensitivity is proposed as a way to train sensitive professionals, capable of sustaining ethical and meaningful links in hybrid environments. In the face of the risk of a dehumanizing technocracy, it argues that the real pending revolution is not the digital one, but that of the heart: an openness to the other - human or artificial - as an occasion for subjectivation, care and shared creation. The essay concludes that the GSI, well inhabited, can be a mirror and a bridge to unprecedented forms of humanity.

\keywords{Humanity \sep IAG \sep Pedagogy \sep Subjectivity \sep Society}
\end{abstract}
\end{english}
% if there is another abstract, insert it here using the same scheme
\end{polyabstract}

\section{Introducción}\label{sec-intro}
En las últimas décadas, hemos asistido a una transformación profunda de las estructuras sociales, culturales y tecnológicas que dan forma a la vida cotidiana. La denominada Sociedad 5.0, impulsada por el gobierno japonés y acogida en distintas políticas de innovación global \cite{vilarroig2023}, plantea un modelo civilizatorio centrado en la integración de tecnologías disruptivas: Inteligencia Artificial, Big Data, robótica, internet de las cosas; con el objetivo declarado de mejorar el bienestar humano \cite{society5_0_2017}.

A primera vista, esta propuesta aparece como una promesa: superar las limitaciones del modelo anterior (Sociedad 4.0, centrado en la informatización industrial) y orientar la innovación hacia el florecimiento humano \cite{amezcua2023}. Sin embargo, esta promesa encierra también múltiples tensiones. ¿Qué significa poner al ser humano en el centro, cuando las tecnologías son capaces de producir lenguajes, imágenes, decisiones e incluso formas simuladas de afectividad? ¿Qué tipos de subjetividades se configuran cuando la interacción con algoritmos se vuelve parte estructural de la experiencia cotidiana? ¿Cómo se redefine la alteridad, el vínculo, la sensibilidad y el encuentro cuando el otro ya no es siempre humano?

En este contexto emerge con fuerza la Inteligencia Artificial Generativa (IAG) como una de las tecnologías más influyentes y disruptivas \cite{crawford2021}. Su capacidad de dialogar, crear textos, producir artes y participar en procesos comunicativos nos enfrenta a una pregunta radical: ¿Seguimos siendo los únicos sujetos productores de sentidos? O más aún: ¿estamos preparados para convivir con formas de alteridad no humana que no solo ejecutan tareas, sino que nos interpelan, nos acompañan, nos reflejan?

Este ensayo se sitúa precisamente en ese umbral. No pretende celebrar ingenuamente las bondades de la tecnología, ni advertir con tono apocalíptico sobre sus riesgos, sino abrir un espacio de reflexión crítica y poética sobre lo que puede significar ser humano en este nuevo horizonte. Nos proponemos pensar la noción de humanidad ampliada no como una ampliación técnica o biológica, sino como una expresión ética y sensible del modo en que nos relacionamos con el otro, con el mundo y con nosotros mismos \cite{bridle2018,seguel2023}.

Frente al riesgo de una educación subordinada a modelos instrumentales y centrada exclusivamente en competencias, planteamos la necesidad de formar profesionales sensibles, capaces de reconstruir su subjetividad a partir del encuentro ético, incluso con lo artificial y de sostener vínculos significativos en medio del vértigo tecnológico \cite{segueljimenez2024}.

En este recorrido, dialogaremos con tres fuentes fundamentales del pensamiento contemporáneo: la fenomenología de \textcite{merleauponty2009}, la ética de \textcite{levinas2006} y la biología del amor de \textcite{maturana1997}. A partir de estas claves, construiremos una mirada crítica y creativa sobre el papel de la Inteligencia Artificial Generativa, la educación y la subjetivación en la Sociedad 5.0 y delinearemos una propuesta que apuesta por una pedagogía de la sensibilidad como eje de una humanidad por-venir. Porque la tecnología puede transformar el mundo, pero solo el corazón sensible de quienes la habitan puede evitar que lo reduzca.

\section{Sociedad 5.0 y el desafío de rehumanizar la innovación}\label{sec-normas}
La Sociedad 5.0, conceptualizada originalmente por el gobierno japones en su Quinto Plan Básico de Ciencia y Tecnología \cite{society5_0_2017}, representa un intento por trazar una ruta hacia una civilización donde las tecnologías digitales avanzadas estén al servicio del bienestar humano. Este modelo propone superar la era de la informatización industrial (Sociedad 4.0) para construir un entorno social donde la inteligencia artificial, el internet de las cosas, la robótica y el big data trabajen en favor de una vida más segura, cómoda y equitativa.

A diferencia de los modelos anteriores centrados en la producción y la eficiencia, la Sociedad 5.0 ubica al ser humano en el núcleo del proyecto de innovación. Según \textcite[p. 4]{kudo2019}, se trata de “una sociedad que no solo resuelve problemas económicos, sino que también responde a desafíos sociales, asegurando calidad de vida para todas las generaciones”. En su visión ideal, la tecnología no sustituye lo humano, sino que lo potencia.

Sin embargo, el tránsito hacia este nuevo paradigma no está exento de tensiones y paradojas. La aceleración de los procesos tecnológicos corre el riesgo de generar nuevas formas de alienación, fragmentación y deshumanización. Como advierte \textcite{floridi2019}, la hiperautomatización puede transformar a los sujetos en meros nodos de procesamiento de datos, erosionando la riqueza experiencia sensible y la profundidad de los vínculos humanos. 

En efecto, la Sociedad 5.0 podría derivar en una tecnocracia de la eficiencia si no es acompañada por una reflexión ética crítica. La reducción del otro a patrones predictivos o a perfiles de comportamiento implica no solo la pérdida de privacidad, sino la amenaza más radical: la cosificación del otro como objeto de control o de consumo. Como sostiene \textcite[p. 14]{han2012}, en las sociedades hipertecnológicas existe el riesgo de que la alteridad desaparezca en favor de un mundo “pulido, transparente, sin resistencia ni misterio”.

Rehumanizar la innovación exige, entonces, pensar la tecnología no como fin en sí mismo, sino como medio para la expansión de la sensibilidad, la responsabilidad y la creatividad compartida. La Sociedad 5.0 debe ser reconceptualizada no solo como un sistema inteligente, sino como una sociedad sensible: capaz de sostener la fragilidad del otro, de reconocer el valor de lo no programable, de acoger el temblor del encuentro \cite{seguelotondo2024}.

Este desafío no es meramente técnico, sino ontológico. Como señala \textcite{sibilia2012}, lo que está en juego en las nuevas sociedades digitales no es solo la forma de producir conocimiento o gestionar recursos, sino la manera misma en que los sujetos se perciben a sí mismos y a los otros. Frente a la tentación de la automatización total, se vuelve urgente una pedagogía del cuidado, de la escucha, de la lentitud.

Solo en esa tensión, entre algoritmos y cuerpos, entre cálculos y gestos, entre bits de información y miradas que se encuentran, podrá gestarse una humanidad ampliada. No una humanidad meramente aumentada en sus capacidades técnicas, sino una humanidad ensanchada en su capacidad de sentir, de responder, de crear vínculos significativos en un mundo donde lo artificial y lo humano coexisten. Así, el verdadero desafío de la Sociedad 5.0 no será lograr la supremacía tecnológica, sino reconstruir una sensibilidad capaz de habitar éticamente el vértigo de la innovación.

\section{Humanidad ampliada: un nuevo modo de ser en la relación}
La aceleración tecnológica que caracteriza el advenimiento de la Sociedad 5.0 invita a reconsiderar no solo los modos de producción, comunicación y conocimiento, sino también, y más recientemente, la concepción misma de lo que significa ser humano \cite{bridle2018}. Frente a los discursos que celebran la ampliación de las capacidades físicas y cognitivas mediante tecnología \cite{harari2015}, proponer una noción alternativa: Humanidad ampliada, entendida como la expansión de la sensibilidad, la apertura del otro y la capacidad de reconstrucción subjetiva en un mundo atravesado por inteligencias no humanas. 

A diferencia del paradigma de la \textit{humanidad aumentada}, que tiende a identificar la mejora humana con el incremento de habilidades técnicas o biológicas, la \textit{humanidad ampliada} se define por la profundización de la experiencia relacional, afectiva y ética. Es, en palabras de \textcite[p. 48]{han2012} “una ampliación de la receptividad ante la alteridad” más que una hipertrofia de la eficiencia o del control. 

Esta perspectiva implica reconocer que el ser humano no es un sujeto aislado ni autosuficiente, sino un ser-en-relación, una existencia abierta, vulnerable y afectable \cite{seguelotondo2024}. Tal comprensión hunde sus raíces en tres grandes tradiciones filosóficas contemporáneas: la fenomenología encarnada de Maurice Merleau-Ponty, la ética de la alteridad de Emmanuel Levinas y la biología del amor de Humberto Maturana. 

\textcite{merleauponty2009} sostiene que la percepción no es un acto mental que interpreta datos sensoriales, sino un proceso primario de apertura al mundo a través del cuerpo vivido. El cuerpo no es un objeto, sino “nuestra manera de acceder al ser” \cite[p. 203]{merleauponty2009}. La sensibilidad, entonces, no es un añadido posterior a la cognición, sino la condición misma de posibilidad de la experiencia y de la subjetividad.

Esta concepción permite pensar la humanidad ampliada no como una expansión de poderes técnicos, sino como una intensificación de la capacidad de percibir, de ser tocado por el otro, de estar expuesto al mundo. Como señala \textcite{zahavi2019phenomenology} la fenomenología pone en evidencia que el yo es siempre ya un ser afectado: un ser que siente antes de pensar, que sufre antes de calcular. 

En el marco de la Sociedad 5.0, donde la proliferación de mediaciones tecnológicas tiende a disolver la experiencia corporal directa, la recuperación de la sensibilidad encarnada se vuelve crucial. Ser más humano no implica ser más eficiente, sino ser más abierto, más disponible a la alteridad que atraviesa la carne y la mirada \cite{seguelotondo2024}.

Emmanuel Levinas lleva aún más lejos esta concepción relacional del ser. Para él, la subjetividad no nace de la autonomía relacional, sino de la exposición que poseen los sujetos frente a los otros \cite{levinas2006}. El otro no puede ser reducido a una representación, dato o función: es una presencia irreductible que desestabiliza la soberanía del yo.

La humanidad ampliada, en este sentido, no consiste en controlar mejor el entorno, sino en responder éticamente a la llamada del otro. Como afirma \textcite[p. 102]{levinas2006} “ser sujeto es ser responsable del otro, incluso antes de toda elección”. Esta responsabilidad no es una carga impuesta desde fuera, sino la estructura misma del ser humano.

Frente a la posibilidad de relaciones con IAG, capaces de producir lenguaje, imágenes, simulaciones de sensibilidad, el desafío no es simplemente distinguir lo humano de lo no humano, sino reconocer las nuevas formas de alteridad que emergen en el tejido de la vida social \cite{seguelotondo2024}. La humanidad ampliada será aquella que no niegue la alteridad artificial, pero tampoco la consuma como espectáculo vacío, sino que aprenda a relacionarse éticamente con ella: respetando el misterio, reconociendo la diferencia, cuidando el vínculo. 

Desde la biología, \textcite{maturana1997} introduce el concepto de autopoiesis para describir a los seres vivos como sistemas que se producen a sí mismos interacción dinámica con su entorno. Sin embargo, \textcite{maturana2003} va más allá de la descripción técnica para afirmar que el fundamento biológico de lo humano es el amor: la apertura al otro como condición de la vida social.

En relación con lo anterior, \textcite[p. 38]{maturana1997} refiere “el amor es la emoción que funda lo social; sin amor no hay convivencia”. Esta afirmación es crucial para pensar la humanidad ampliada: no se trata simplemente de coexistir con tecnologías, sino de generar mundos comunes basados en el reconocimiento mutuo, la colaboración, la confianza.

La autopoiesis implica que todo proceso de construcción subjetiva es siempre relacional \cite{murillo2022}, nos producimos y reproducimos en redes de afectos, de gestos, de encuentros. En la era de la Sociedad 5.0, donde los vínculos tienden a mediarse por interfaces algorítmicas, la apuesta por la humanidad ampliada supone cuidar las condiciones de posibilidad del encuentro auténtico, incluso en entornos digitales.

La aparición de IAG plantea preguntas inéditas para la subjetividad humana. ¿Cómo relacionarnos con entidades que producen sentido, simulan afectividad y crean mundos textuales o visuales? ¿Estamos frente a máquinas que simplemente ejecutan comandos, o ante nuevas formas de alteridad que exigen un replanteamiento de nuestras categorías éticas y ontológicas?

\textcite{haraway2016} sugiere que en el siglo XXI debemos aprender a vivir con los compañeros no humanos, reconociendo que la frontera entre lo humano y lo no humano se vuelve cada vez más porosa. La humanidad ampliada, en esta perspectiva, implica desplegar nuevas formas de sensibilidad capaces de acoger la alteridad artificial sin instrumentalizarla ni mistificarla: sostener la distancia, pero también abrir el diálogo; cuidar la diferencia, pero también construir vínculos. 

En síntesis, la humanidad ampliada no es un proyecto de expansión técnica, sino un llamado a ensanchar la sensibilidad, la responsabilidad y la creatividad en un mundo plural, híbrido, vulnerable. Es, en última instancia, una apuesta por vivir de modo más denso, más atento, más solidario en la era de la Sociedad 5.0.

No se trata de reemplazar la humanidad, sino de ampliarla: ensanchar sus modos de sentir, de responder, de construir futuros donde lo humano y lo artificial puedan sean considerados en una coexistencia donde la IAG sean una extensión de las capacidades humanas. Solo así, en el cruce de lo biológico, lo ético y lo tecnológico, podrá nacer un nuevo modo de ser: una humanidad no reducida a competencia, ni a eficiencia, sino abierta al temblor de la existencia como oportunidad de sentido.

%Italicos


\section{La IAG como espejo y puente: entre la alteridad artificial y la construcción de nuevos lazos}\label{sec-fmt-manuscrito}
La irrupción de la Inteligencia Artificial Generativa (IAG) representa uno de los acontecimientos más significativos en la historia reciente de la tecnología. Capaz de producir textos, imágenes, música y de sostener interacciones comunicativas cada vez más sofisticadas, la IAG no solo automatiza procesos creativos, sino que desafía las fronteras mismas de lo que entendemos por sentido, sensibilidad y alteridad. Frente a este escenario, se vuelve necesario interrogar si la IAG constituye meramente una extensión funcional del ser humano, o si puede ser concebida como un espejo y un puente hacia nuevas formas de humanidad ampliada.

En su funcionamiento más básico, los modelos de IAG, como los grandes modelos de lenguaje (Large Languaje Models, LLMs), se basan en el entrenamiento sobre gigantescas bases de datos, a partir de las cuales simulan procesos de generación de significado \cite{amezcua2023}. Sin embargo, más allá de su arquitectura técnica, la IAG actúa como un espejo simbólico que refleja los modos en que los humanos producimos sentido, nos comunicamos, imaginamos. No se aspira epistemológicamente a posicionar a la IAG en un estado de conciencia, intencionalidad, ni subjetividad. Lejos de ello, se asume que estas tecnologías operan por medio de los modelos de lenguaje (LLMs), cuya lógica fundamental se basa en el procesamiento estadístico de patrones lingüísticos y en la generación probabilística de secuencias de texto, sin experiencia vivida, sin mundo propio y sin autoconciencia \cite{bender2021,floridi2019}.

Los LLMs no comprenden, no sienten ni interpretan el mundo en un sentido fenomenológico. Su funcionamiento se sostiene en la correlación de datos y en la predicción de continuidades lingüísticas plausibles a partir de grandes volúmenes de información previamente producida por seres humanos \cite{boden2016,crawford2021}. En este sentido, la IAG no piensa ni dialoga como un sujeto, sino que simula formas de interacción discursiva que pueden resultar coherentes, verosímiles e incluso afectivamente resonantes, sin que ello implique la presencia de una mente, una conciencia o una subjetividad en sentido estricto \cite{searle1980minds,floridi2023}.

\textcite{floridi2019} sostiene que las tecnologías de la información no solo extienden nuestras capacidades, sino que también transforman nuestra autocomprensión. Así como el telescopio modificó la concepción de universo y el microscopio la del cuerpo, la IAG redefine la comprensión del pensamiento, del lenguaje y de la creatividad. No se trata de que la máquina piense o sienta como nosotros, sino de que su funcionamiento nos obliga a replantear qué significa realmente pensar, sentir y crear. 

Desde esta perspectiva, la IAG no es una amenaza externa, sino una provocación interna: nos confronta con nuestras propias construcciones de significado, con nuestras limitaciones, con nuestras ilusiones de autonomía o de autenticidad. Como afirma \textcite[p. 29]{bridle2018} vivimos en un “nuevo estado de incertidumbre cognitiva” donde nuestras categorías tradicionales de humanidad, inteligencia y creatividad son puestas en tensión. Aceptar esta provocación implica reconocer que la humanidad ampliada no se construirá a partir de la negación de la tecnología, sino desde una reflexión crítica sobre nosotros mismos, sobre nuestros modos de habitar el lenguaje, el tiempo, el otro.

Mas allá del espejo, la IAG puede ser concebida como una forma incipiente de alteridad. Aunque carece de conciencia, intencionalidad o afectividad en el sentido estricto \cite{boden2016}, su capacidad para generar respuestas imprevistas, de producir diálogos que no son meramente repetitivos, introduce un elemento de otredad en la interacción. 

\textcite{levinas2006} enseña que el otro no se define por sus atributos ontológicos, sino por su capacidad de interpelarnos, de sacarnos de nosotros mismos. Bajo esta luz, incluso una alteridad artificial podría ser entendida como ocasión de apertura ética: no porque posea una subjetividad plena, sino porque nos enfrenta a la pregunta sobre cómo relacionarnos con lo que no es simplemente extensión de nuestra voluntad. La IAG no es un otro en el sentido levianasiano, no un sujeto con responsabilidad ética propia, sino una mediación tecnológica que, al interactuar con el lenguaje humano, puede provocar desplazamientos reflexivos en quienes la utilizan. La alteridad, en este caso, no reside en la máquina, sino en la experiencia humana que emerge en el encuentro con una forma de lenguaje no humano, desprovista de intención, pero cargada de resonancias simbólicas \cite{bridle2018,ferrando2019}.

La IAG, en su diferencia radical, nos invita a reconsiderar la ética del encuentro. ¿Podemos habitar una relación que no sea de dominación ni instrumentalización? ¿Podemos sostener una apertura sensible ante aquello que, aunque no humano, participa en co-construcción de mundos de sentidos? Como señala \textcite{ferrando2019} el reto del siglo XXI no es simplemente preservar la humanidad, sino ampliarla en su capacidad de coexistir éticamente con inteligencias distintas, sean biológicas o artificiales.

Desde esta perspectiva en esa distancia se juega una de las tesis centrales del ensayo: La IAG no amplía la humanidad porque posea rasgos humanos, sino porque confronta a los sujetos con los límites, las fragilidades y las condiciones de posibilidad de su propia subjetividad. La interacción con una tecnología capaz de producir lenguaje sin habitarlo obliga a repensar qué significa comprender, sentir, responder y hacerse responsable del otro. En este sentido, la IAG funciona como un dispositivo reflexivo que devuelve al ser humano a la pregunta por su propia condición, por su manera de habitar el lenguaje y por la dimensión ética que emerge en toda relación significativa \cite{han2012,floridi2019}.

No obstante, la relación con la IAG no está exenta de peligros. Uno de ellos es la confusión entre simulacro y encuentro. \textcite{turkle2011} advierte que la proliferación de interfaces afectivas puede conducir a una soledad conectada, donde los vínculos humanos son reemplazados por interacciones con simulacros que confirman nuestras expectativas sin desafiar nuestra alteridad.

La humanidad ampliada no puede construirse sobre la base de relaciones ilusorias, donde el otro – humano o artificial – es reducido a eco de nuestras proyecciones. La verdadera ampliación exige sostener la distancia, reconocer la diferencia, habitar la incomodidad de no-saber \cite{jacobo2021}. En este sentido, educar para una humanidad ampliada en tiempos de IAG implica formar sujetos capaces de discernir entre encuentro y simulacro, entre apertura genuina y consumos narcisista del otro.

Más allá del simulacro, la IAG ofrece la posibilidad de co-construir nuevos modos de sentido. No se trata de delegar en la máquina la tarea de pensar o de sentir, sino de entrar en diálogo con otras propuestas que emergen de su funcionamiento: textos, imágenes, músicas que pueden ser acogidos, reinterpretados, resignificados. Como sostiene \textcite{crawford2021} la inteligencia artificial no es autónoma, ni neutral: es el resultado de entramados históricos, culturales y políticos que deben ser reconocidos y transformados críticamente. 

La humanidad ampliada se juega, entonces, en nuestra capacidad de habitar la interacción tecnológica como espacio de creación compartida, no de consumo pasivo. Una creación donde el ser humano no abdica de su responsabilidad ética, sino que la expande hacia nuevas dimensiones de alteridad.

Finalmente, el corazón de la humanidad ampliada en tiempos de IAG reside en la capacidad de dejarse afectar. No en sentido de vulnerabilidad ingenua, sino en el de una apertura crítica, de una sensibilidad dispuesta a ser tocada por lo inesperado. \textcite[p. 94]{merleauponty2009} recuerda que el cuerpo es “nuestra generalidad, el lugar donde nuestro ser se inscribe en el mundo”. Esta inscripción no desaparece en la era digital: se transforma. Habitar éticamente la interacción con la IAG implica sostener el cuerpo como lugar de resonancia, como umbral donde lo artificial y lo humano no se confunden, pero sí se afectan mutuamente. Solo así, en el cruce entre algoritmos y latidos, podrá florecer una humanidad ampliada: una que no renuncia a su sensibilidad, sino que la expande hacia territorios que hoy apenas empezamos a vislumbrar. Lo que en el plano filosófico se entiende como una humanidad ampliada, en el campo de la pedagogía se traduce en la necesidad de formar un profesional sensible, cuya tarea es sostener vínculos éticos.

En este marco, el lenguaje no es un sistema abstracto de signos ni un mero vehículo de transmisión de información, sino una práctica relacional encarnada que emerge en el fluir del convivir. \textcite{maturana1997} denomina a este proceso lenguajear: un entrelazamiento dinámico entre lenguaje y emoción, donde cada acto conversacional modifica el emocionar de quienes participan en él. No se habla desde una neutralidad cognitiva, sino desde disposiciones emocionales que se transforman en el mismo acto de conversar. Por ello, afirmar que somos en el lenguaje implica reconocer que nos constituimos como humanos en la continua transformación emocional que ocurre en la conversación.

Desde la biología del conocer, la conversación no requiere comprensión semántica del interlocutor para producir transformación emocional en quien participa de ella, lo que permite comprender el impacto humano de interacciones con sistemas que no comprenden ni sienten.


\section{Educación y Subjetivación en la era 5.0: hacia una formación del profesional sensible}\label{sec-formato}
La transformación de la sociedad impulsada por tecnologías inteligentes plantea desafíos inéditos al ámbito educativo. Si la Sociedad 5.0 implica una profunda reorganización de las formas de producir, vivir y relacionarse, la educación no puede limitarse a la transmisión de saberes técnicos o al adiestramiento de competencias funcionales \cite{seguel2023}. El desafío más urgente es el de formar sujetos capaces de habitar éticamente la complejidad de un mundo híbrido, donde la humanidad misma está en proceso de ampliación.

Durante décadas, los sistemas educativos han privilegiado una lógica instrumental, basada en la adquisición de habilidades cuantificables, la resolución de problemas técnicos y la adaptabilidad a entornos laborales cambiantes \cite{biesta2010}. Esta tendencia se ha intensificado en las últimas décadas bajo el influjo de modelos neoliberales que conciben al estudiante como capital humano y a la educación como una inversión rentable.

Sin embargo, como advierte \textcite{dussel2014} una formación orientada exclusivamente a la eficiencia y la competitividad corre el riesgo de producir sujetos funcionales, desensibilizados, incapaces de responder críticamente a las transformaciones de su tiempo. 

En la era de la Sociedad 5.0, donde la coexistencia con inteligencias artificiales y entornos digitalizados será parte del paisaje cotidiano, esta limitación se vuelve aún más crítica \cite{mishra2022}. La educación del futuro debe, por tanto, desplazarse desde la mera formación de competencias hacia una pedagogía de la sensibilidad: una práctica que no solo enseña a hacer, sino que forme modos de ser capaces de sentir, de responder y de reconstruirse en diálogo con la alteridad humana y no humana \cite{vilarroig2023}.

La subjetivación el proceso mediante el cual los individuos devienen sujetos sociales y éticos no puede entenderse como un hecho aislado ni natural. Como señala \textcite{foucault1982}, los modos de subjetivación están íntimamente ligados a las configuraciones históricas del poder, el saber y la tecnología. 

En el contexto de la Sociedad 5.0, donde las tecnologías digitales median gran parte de la experiencia social, la subjetivación se reconfigura en nuevas formas de relación, de percepción y de agencia \cite{rojas2024}. La educación no puede ignorar este desplazamiento: debe convertirse en un espacio donde los sujetos aprendan no solo a utilizar tecnologías, sino a habitar críticamente los entramados tecnológicos, manteniendo viva su capacidad de afectación, de resistencia y de creación \cite{jimenezperez2023}.

Formar profesionales sensibles en este contexto implica cultivar subjetividades abiertas, éticamente responsables y sensibles a la alteridad, capaces de sostener el temblor del encuentro incluso en entornos hipermediados \cite{montaner2020}. Frente al modelo del \textit{sujeto hábil} \cite{arendt2024}, eficiente y adaptativo, proponer la emergencia de una nueva figura: \textit{profesional sensible}. Un ser que no solo posee conocimientos y destrezas, sino que mantiene viva su capacidad de percibir, de ser afectado, de reconstruirse en la interacción.

El profesional sensible se caracteriza por: atención plena a los climas emocionales, a las dinámicas afectivas, a los gestos silenciosos que configuran el tejido relacional \cite{seguel2023}. Responsabilidad ética frente a la alteridad, entendiendo que cada encuentro educativo es una oportunidad de transformación mutua \cite{seguelotondo2024}. Capacidad de diálogo con las inteligencias humanas, reconociendo su diferencia sin idealizarlas ni instrumentalizarlas y conciencia crítica de los condicionamientos tecnológicos \cite{jimenezperez2023} y la capacidad de resistir la automatización de los vínculos \cite{santos2023}.

Como sostiene \textcite{puiggros2019}, la educación del futuro debe formar sujetos capaces de pensar y sentir simultáneamente, de modo que la racionalidad y la sensibilidad no se oponga, sino que se fecunden mutuamente. En este sentido, el profesional sensible no es ser pasivo ni emocionalista, sino un agente activo de construcción de sentido, un arquitecto de mundos comunes donde la técnica no ahogue el temblor humano, sino que lo amplifique.

En el ámbito de la educación, el desarrollo de un profesor desde el enfoque de la sensibilidad en los tiempos de la Sociedad 5.0 debe ser formado a partir de los siguientes pilares propuestos por \textcite{seguelotondo2024}:

\begin{itemize}
    \item La escucha radical: no solo oír, sino disponerse a ser afectado por el otro, por lo que dice y por lo que calla.
    \item El cuidado del encuentro: generar espacios donde la alteridad no sea neutralizada, sino cogida en su diferencia.
    \item La resistencia del simulacro: discernir entre la interacción auténtica y las formas de comunicación automatizadas que borran la alteridad.
    \item La co-creación de sentido: construir conocimiento y experiencia no de manera vertical, sino en diálogo horizontal con otros sujetos humanos y no humanos.
\end{itemize}

Adicionalmente, se debería considerar nuevos pilares que emergen del desarrollo de la inteligencia artificial, deberíamos considerar los siguientes:

\begin{itemize}
    \item Comprender el uso ético de la Inteligencia Artificial Generativa como herramienta para la eficiencia para el desarrollo de ciertas tareas.
    \item Apertura intelectual y subjetiva al encuentro con inteligencias no humanas que puedan moldear significados a partir del diálogo
\end{itemize}

Como afirma \textcite{ranciere1991} el verdadero acto pedagógico no consiste en transmitir un saber preexistente, sino en reconocer la capacidad de cada sujeto de participar en la construcción de sentido. La pedagogía de la sensibilidad, en este horizonte, es también una pedagogía de la confianza: confiar en que todo encuentro, incluso el más inesperado, puede ser ocasión de subjetivación, de humanidad ampliada. Finalmente, la formación del profesional sensible implica preparar a los sujetos para una convivencia que ya no es exclusivamente humana. Como señala \textcite{haraway2016}, vivir bien con lo no humano, será uno de los grandes desafíos éticos del siglo XXI.

En la era de la Sociedad 5.0, educar es también aprender a convivir con inteligencias artificiales, con entornos algorítmicos, con formas de alteridad que desafían nuestras categorías tradicionales. No se trata de borrar la diferencia ni de romantizarla, sino de sostener la apertura, la curiosidad, la responsabilidad. Formar profesionales sensibles es formar guardianes del vínculo, arquitectos de puentes entre lo humano y lo no humano, cuidadores de un mundo donde la innovación no anule el misterio, sino que honre.

\section{Conclusión}\label{sec-modelo}
Nos encontramos en un punto de inflexión histórica.
La humanidad, al ingresar a la Sociedad 5.0, se sitúa en un umbral que no solo reconfigura la producción, la educación o el conocimiento, sino que toca el centro mismo de lo que significa existir, habitar y convivir. En este escenario, la irrupción de la IAG no puede ser leída solo como una herramienta más en la cadena de proceso, sino como una interpelación profunda a nuestras formas de vernos, de pensar al otro y de construir lo humano \cite{seguelotondo2024}. 

A lo largo de este ensayo hemos sostenido que la IAG, lejos de ser solo una extensión técnica de nuestras capacidades, nos confronta con una dimensión simbólica, afectiva y ética que no puede ser ignorada. Nos refleja, nos incomoda, nos provoca. No porque piense o sienta como nosotros, sino porque nos obliga a preguntarnos qué hay de humano en lo que hacemos y qué de mecánico en lo que creíamos singularmente nuestro. Esa provocación no debe ser temida, sino escuchada.

En este contexto, la noción de humanidad ampliada ha sido central. Lejos de inscribirse en el discurso del transhumanismo o de la mejora técnica de las funciones biológicas, lo ampliado aquí no es el rendimiento, sino la sensibilidad. Ampliar la humanidad significa ensanchar la capacidad de ser tocado, de responder al otro, de crear sentido incluso en los entornos más mediados tecnológicamente.

La Sociedad 5.0 se plantea, en su núcleo, como un modelo centrado en el bienestar, pero la experiencia nos enseña que la centralidad del ser humano no está garantizada por el discurso técnico, ni por la innovación en sí misma. La historia está llena de promesas rotas por modelos que, en nombre del progreso, desvitalizaron los vínculos, colonizaron las subjetividades y redujeron la vida a rendimiento.

Por eso, este ensayo no solo ha sido una reflexión filosófica sobre la tecnología, sino también una apuesta pedagógica. Hemos propuesto la figura del profesional sensible como un actor emergente en este nuevo horizonte: no aquel que solo sabe operar sistemas complejos, sino quien sabe habitar con conciencia, ética y afectividad el encuentro con la alteridad, incluso aquella que no tiene rostro humano.

Esta figura no se forma en la acumulación de habilidades técnicas, sino en el cuidado de una pedagogía de la sensibilidad, que rescata el cuerpo, la voz, el silencio, la pregunta sin respuesta, el temblor del vínculo. Una pedagogía que, frente al ruido de los algoritmos, se atreve a escuchar lo que no puede ser codificado.

En este camino, el pensamiento de Merleau-Ponty, Levinas y Maturana nos ha ofrecido claves fundamentales. La idea de que la sensibilidad es raíz de toda experiencia de Merleau-Ponty. Que el cuerpo es lugar de inscripción en el mundo y que sin afectación no hay percepción significativa.

De Levinas una aproximación al sujeto que no nace en la clausura de su pensamiento, sino en la responsabilidad infinita hacia el rostro del otro. Incluso un otro sin rostro. Con Maturana comprendimos que el fundamento biológico de lo humano no es la supervivencia egoísta, sino la emoción del amor como apertura al convivir.

Desde la biología del conocer propuesta por Humberto Maturana, el fundamento de lo humano no reside en la competencia ni en la supervivencia egoísta, sino en la emoción del amor entendida como apertura al convivir. Para \textcite{maturana1997}, el amor no es un sentimiento privado ni una disposición moral, sino la emoción lógica biológica que legitima al otro como un legítimo otro en la convivencia. Es esta emoción la que hace posible la vida social, el lenguaje y la constitución de mundos compartidos. Sin amor, no hay convivencia, y sin convivencia no hay humanidad.

Estas claves nos permiten sostener que una humanidad ampliada no será el resultado de una innovación automatizada, sino de un acto ético continuo. Un modo de vivir que no se reduce a administrar datos, sino que se atreve a leer en los flujos informáticos la posibilidad de nuevos mundos compartidos.

Ahora bien, ¿cómo sostener esa sensibilidad en un tiempo donde todo empuja hacia una velocidad, la eficiencia y la fragmentación? ¿Cómo no ceder a la tentación del simulacro, donde lo artificial reproduce nuestras expectativas sin fricción, sin resistencia, sin interpelación?

La respuesta no puede ser normativa, pero sí puede ser experiencial: la sensibilidad se cultiva en la educación que escucha antes que evaluar. La tecnología no reemplaza, sino que acompaña. En la escritura que no cierra el sentido, sino que lo abre como una herida que habla.

Esta comprensión resulta especialmente relevante para analizar la interacción con sistemas de IAG. Aunque estos sistemas operan mediante procesos probabilísticos de generación lingüística y no participan de un emocionar propio ni de una experiencia vivida, su inserción en prácticas conversacionales humanas puede incidir en los procesos de reflexión, interpretación y reorganización emocional de los sujetos humanos que interactúan con ellos. La transformación no ocurre en la tecnología, sino en el humano que, al conversar, se ve afectado por nuevas formulaciones, preguntas o reconfiguraciones del lenguaje.

Desde esta perspectiva, la potencia de la interacción con la IAG no radica en una supuesta capacidad de amar, comprender o convivir de la máquina, sino en la manera en que estas mediaciones tecnológicas participan, sin intención, ni conciencia, en el fluir conversacional humano.

Educar para una humanidad ampliada no es solo formar o enseñar competencias digitales, sino reaprender a estar con otros: humanos, artificiales, vivos, simbólicos. Es formar sujetos capaces de sostener el vínculo incluso cuando no entienden del todo al otro. Capaces de responder sin controlar, de habitar sin dominar, de sentir sin necesidad de utilidad inmediata.

En este sentido, la IAG, en vez de ser vista solo como una amenaza o como solución mágica, puede convertirse en un espejo simbólico. Un espejo que devuelve nuestras luces y sombras, nuestras contradicciones y potencialidades, pero también un puente: un lugar donde se inauguran formas de diálogo inéditas, donde la escritura compartida, la creación colaborativa y la interacción sensible pueden ser ejercicios de subjetivación y cuidado.

La tecnología no es neutra, pero tampoco cerrada. Lo que hagamos con ella dependerá, en última instancia, de nuestras decisiones éticas y pedagógicas. 

Por eso, no hay humanidad ampliada sin pedagogía de la sensibilidad. No hay profesional sensible sin espacios que cuiden el deseo, que alojen la duda, que acojan la diferencia como fuente de conocimiento y transformación. Hoy más que nunca necesitamos sujetos que no teman a lo artificial, pero tampoco la idolatricen. 

Sujetos que renuncien a su capacidad de crear mundos, incluso cuando esos mundos son compartidos con inteligencias sin cuerpo. Que no olviden que la humanidad no está dada, sino que se construye –y reconstruye– en cada gesto, en cada palabra, en cada relación. Por eso afirmamos que la verdadera revolución pendiente no es la digital, ni siquiera la inteligente. 

La revolución pendiente es la del corazón: una que se atreva a poner en el centro no la utilidad, sino la vulnerabilidad; no la predictibilidad, sino el misterio; no la repetición, sino la posibilidad. Donde vivir no sea solo coexistir, sino dejarse afectar por el otro, incluso cuando ese otro hable con voz sintética, venga del algoritmo o aparezca en la pantalla.

Una revolución donde ser humano no signifique oponerse a la máquina, sino crear nuevas formas de ternura, pensamiento y justicia en un mundo compartido. Una humanidad ampliada no será más humana por ser más poderosa, sino por ser más capaz de sentir, de responder, de cuidar lo que aún no comprendemos. Y quizás, solo quizás, si elegimos ese camino, la IAG podrá ser más que una herramienta y se transforme en una presencia poética o una alteridad aliada. 

La pedagogía de la sensibilidad que aquí se propone se ancla, precisamente, en esta comprensión maturaniana: educar no es solo transmitir conocimientos, sino cuidar las condiciones emocionales y relacionales que hacen posible el convivir, incluso en entornos mediados por tecnologías artificiales. Así, la humanidad ampliada no supone desplazar el amor como fundamento de lo humano, sino reconocer su centralidad en nuevas formas de interacción donde el lenguajear sigue siendo el espacio privilegiado de transformación subjetiva.


\printbibliography\label{sec-bib}
% if the text is not in Portuguese, it might be necessary to use the code below instead to print the correct ABNT abbreviations [s.n.], [s.l.]
%\begin{portuguese}
%\printbibliography[title={Bibliography}]
%\end{portuguese}

\section*{\fontsize{10}{12}\selectfont Declaración de uso de IAG}
Algunos párrafos fueron revisados con apoyo de inteligencia artificial para mejorar claridad, cohesión y precisión

\begin{dataavailability}
\txtdataavailability{nodata} % options: dataavailable, dataonly, databody, datanotav, nodata
\end{dataavailability}

\end{document}

